% Every Lie group is a homogeneous space, not conversely.
%
% The geometry is chosen so that every label has room, and every label position is checked
% against the ellipse it must stay inside:
%   outer (green)  centre (0,0),     rx = 5.5, ry = 1.9
%   inner (blue)   centre (0,-0.45), rx = 3.0, ry = 1.0
% A point (x,y) is inside an ellipse iff (x/rx)^2 + ((y-cy)/ry)^2 < 1, and it is the
% *corner* of a label's box that has to satisfy it, not its centre.  The clearances that
% follow from that (all in cm, all measured rather than argued):
%   `Homogeneous Spaces' at y = 1.30, half-width 2.3: the green boundary is 3.94 wide at
%       y = 1.30 and 3.36 at the box top y = 1.50 -- 1.06 cm of slack at the worst point.
%   `Lie Groups' at y = -0.12, half-width 1.35: the blue boundary is 2.86 wide there,
%       2.72 at the box top y = 0.08 -- 1.37 cm of slack.
%   $SO(n)$, $\mathbb{T}^n$ at (-1.45,-1.00) / (1.45,-1.00), half-width 0.55: blue is
%       2.60 wide at y = -1.00 and 2.29 at the box bottom y = -1.18 -- both corners in.
%   $S^n$ at (-4.30,0.05) and $Gr(n,N)$ at (-3.75,-0.95): the green-blue ring runs
%       3.00..5.49 at y = 0.05 and 2.68..4.85 at y = -0.95.
%   $St(n,N)$ at (4.25,-0.45): the ring runs 3.00..5.35 there.
\documentclass[tikz]{standalone}
\usepackage{geometricfigures}

\usepackage{amsmath}
\usepackage{amssymb}
\usepackage{amsfonts}

\begin{document}
\begin{tikzpicture}[
        label_style/.style={font=\large\bfseries},
        example_style/.style={font=\small, inner sep=1pt},
    ]
    % ---- the two regions ----------------------------------------------------
    \draw[draw=mgreen!60!fg, fill=mgreen!20!bg, very thick] (0,0)     ellipse (5.5cm and 1.9cm);
    \draw[draw=mblue!80!fg,  fill=mblue!20!bg,  very thick] (0,-0.45) ellipse (3.0cm and 1.0cm);

    % ---- names --------------------------------------------------------------
    \node[label_style, text=mgreen!45!fg] at (0, 1.30) {Homogeneous Spaces};
    \node[label_style, text=mblue!80!fg]  at (0,-0.12) {Lie Groups};

    % ---- examples: inside the blue ellipse ----------------------------------
    \node[example_style, text=mblue!80!fg] at (-1.45,-1.00) {$SO(n)$};
    \node[example_style, text=mblue!80!fg] at ( 1.45,-1.00) {$\mathbb{T}^n$};

    % ---- examples: in the green ring, outside the blue ellipse --------------
    \node[example_style, text=mgreen!35!fg] at (-4.30, 0.05) {$S^n$};
    \node[example_style, text=mgreen!35!fg] at (-3.75,-0.95) {$Gr(n,N)$};
    \node[example_style, text=mgreen!35!fg] at ( 4.25,-0.45) {$St(n,N)$};
\end{tikzpicture}
\end{document}
